% mode: ini
% Aimed at new_hyph_exceptions (sections 1111-1118): "Not a letter", "Improper \hyphenation will be flushed", \char and \chardef'd letters, char_num's reswitch, a one-letter word (not entered, and its list stays in memory),
% lengths at and past 63, hyphens at the ends and doubled, the repeat that replaces, an upper case letter through \lccode, spaces, macros, a language, and the words then hyphenated by those exceptions, including with a
% font's character codes through \lccode remapping.
\input prelude
\lefthyphenmin=1 \righthyphenmin=1 \defaulthyphenchar=`\- \hyphenpenalty=-10000 \hsize=1pt \tolerance=10000 \parfillskip=0pt \parindent=0pt
\tenrm \hyphenchar\tenrm=`\-
\def\try#1{\setbox0=\vbox{\hskip0pt #1\par}\showbox0 }
\hyphenation{ta-ble}
\try{table}
% Errors, and the word carries on without the bad part.
\hyphenation{ta!ble}
\hyphenation{a!}
\hyphenation{1a-b}
\hyphenation{!}
\hyphenation{a1b-c}
\hyphenation{a\relax b}
\hyphenation{ab$cd}
\hyphenation{ab{cd}ef}
\hyphenation{ab\hskip 1pt cd}
\hyphenation{\undefined}
% \char and a \chardef in a word; a \char of a character with lccode 0.
\hyphenation{\char`a-\char`b c\char`d}
\chardef\cq=`q
\hyphenation{\cq-a\cq}
\hyphenation{\char`!}
\hyphenation{\char 300}
\try{abcd}
\try{qaq}
% One letter: not entered.  Empty.  Only hyphens.
\hyphenation{a}
\hyphenation{a- -a}
\hyphenation{}
\hyphenation{-}
\hyphenation{---}
\hyphenation{ }
\hyphenation{   }
% Hyphens at both ends and doubled; positions out of the word range never set a hyphen.
\hyphenation{-ab -bc- c--d}
\try{ab}
\try{bc}
\try{cd}
% Upper case through \lccode, and the same word in small letters finds it.
\hyphenation{UP-CASE}
\try{upcase}
\try{UPCASE}
\try{UpCaSe}
\uchyph=0
\try{UPCASE}
\try{Upcase}
\uchyph=1
\try{Upcase}
% Repeats: the new positions win, and the count stays.
\hyphenation{re-peat}
\try{repeat}
\hyphenation{rep-eat}
\try{repeat}
\hyphenation{repeat}
\try{repeat}
% Other languages keep their own; the language is part of the word.
\language=1
\hyphenation{re-p-eat}
\try{repeat}
\language=255
\hyphenation{r-epeat}
\try{repeat}
\language=300
\hyphenation{repea-t}
\try{repeat}
\language=0
\try{repeat}
% Sixty-three and past it: the hyphens and letters over the limit are dropped.
\hyphenation{aaaaaaaaaaaaaaaaaaaaaaaaaaaaaaaaaaaaaaaaaaaaaaaaaaaaaaaaaaaaaa-a}
\hyphenation{bbbbbbbbbbbbbbbbbbbbbbbbbbbbbbbbbbbbbbbbbbbbbbbbbbbbbbbbbbbbbbb-bbbbb}
\hyphenation{cccccccccccccccccccccccccccccccccccccccccccccccccccccccccccccccccccc-}
\try{aaaaaaaaaaaaaaaaaaaaaaaaaaaaaaaaaaaaaaaaaaaaaaaaaaaaaaaaaaaaaaa}
\try{bbbbbbbbbbbbbbbbbbbbbbbbbbbbbbbbbbbbbbbbbbbbbbbbbbbbbbbbbbbbbbb}
\try{cccccccccccccccccccccccccccccccccccccccccccccccccccccccccccccccc}
% A macro expanding to two words and a space by \space.  Then one that makes a word by \lowercase.
\def\ws{ma-cro ex-pand}
\hyphenation{\ws}
\try{macro}
\try{expand}
\def\mk{\let\sp= }\mk
\hyphenation{so-me\sp ot-her}
\try{some}
\try{other}
% A missing brace.
\hyphenation x-y}
\hyphenation\relax{y-z}
\let\bg={
\hyphenation\bg z-w}
\try{zw}
% Non-letters to \lccode: a dash as a letter is still a hyphen here.
\lccode`\-=`\-
\hyphenation{d-e}
\try{de}
\lccode`\-=0
% A word needing a font character: a hyphenated word with no such exception and no patterns of the language.
\try{nothing}
\language=7
\try{nothing}
\end
